% Part: first-order-logic
% Chapter: sequent-calculus
% Section: soundness-identity

\documentclass[../../../include/open-logic-section]{subfiles}

\begin{document}

\olfileid{fol}{seq}{sid}

\olsection{Kasahihan mawa \usetoken{S}{identity}}

\begin{prop}
$\Log{LK}$ mawa sekuen wiwitan lan aturan kanggo identitas iku sah.
\end{prop}

\begin{proof}
Sekuen wiwitan kang awujud ${} \Sequent \eq[t][t]$ iku valid, amarga
kanggo saben !!{structure}~$\Struct M$, $\Sat{M}{\eq[t][t]}$. (Cathet
yen kita nganggep term $t$ iku tertutup, yaiku ora ngemot variabel,
mula penetapan variabel ora relevan).

Upamana inferensi pungkasan ing !!a{derivation} iku $=$. Banjur premis
yaiku $\eq[t_1][t_2], \Gamma \Sequent \Delta, !A(t_1)$ lan
kesimpulane yaiku $\eq[t_1][t_2], \Gamma \Sequent \Delta, !A(t_2)$.
Nimbang !!a{structure}~$\Struct M$. Kita kudu nuduhake yen kesimpulane
valid, yaiku yen $\Sat{M}{\eq[t_1][t_2]}$ lan $\Sat{M}{\Gamma}$, mula
salah siji $\Sat{M}{!C}$ kanggo sawetara $!C \in \Delta$ utawa
$\Sat{M}{!A(t_2)}$.

Miturut hipotesis induksi, premis kasebut valid. Iki ateges yen
$\Sat{M}{\eq[t_1][t_2]}$ lan $\Sat{M}{\Gamma}$, salah siji (a) kanggo
sawetara $!C \in \Delta$, $\Sat{M}{!C}$ utawa (b)
$\Sat{M}{!A(t_1)}$. Ing kasus (a) bukti kita rampung. Nimbang kasus
(b). Ayo $s$ dadi penetapan variabel kanthi $s(x) =
\Value{t_1}{M}$. Miturut
\olref[syn][ass]{prop:sentence-sat-true},
$\Sat{M}{!A(t_1)}[s]$. Amarga $\varAssign{s}{s}{x}$, miturut
\olref[syn][ext]{prop:ext-formulas}, $\Sat{M}{!A(x)}[s]$. Amarga
$\Sat{M}{\eq[t_1][t_2]}$, kita duwe $\Value{t_1}{M} =
\Value{t_2}{M}$, mula $s(x) = \Value{t_2}{M}$. Kanthi ngetrapake
\olref[syn][ext]{prop:ext-formulas} maneh, kita uga duwe
$\Sat{M}{!A(t_2)}[s]$. Miturut
\olref[syn][ass]{prop:sentence-sat-true},
$\Sat{M}{!A(t_2)}$.
\end{proof}

\end{document}
